\origpage{43}{N. Robertson, P.D. Seymour / Journal of Combinatorial Theory, Series B 89 (2003) 43--76}
{\centering
{\Large Graph Minors. XVI. Excluding a non-planar graph\par}\medskip

Neil Robertson$^{a,1}$ and P.D. Seymour$^{b,2}$\par
$^a$\textit{Ohio State University, Columbus, OH 43210, USA}\par
$^b$\textit{Telcordia Technologies, 445 South St., Morristown, NJ 07960, USA}\par
Received 4 April 1991\par
}
\begin{quote}
\textbf{Abstract}

This paper contains the cornerstone theorem of the series. We study the structure of graphs with no minor isomorphic to a fixed graph $L$, when $L$ is non-planar. (The case when $L$ is planar was studied in an earlier paper.) We find that every graph with no minor isomorphic to $L$ may be constructed by piecing together in a tree-structure graphs each of which ``almost'' embeds in some surface in which $L$ cannot be embedded.
\end{quote}

\sectionhead{1. Introduction}

Graphs in this paper are finite, undirected graphs which may have loops or multiple edges. A graph $L$ is a \textit{minor} of a graph $G$ if $L$ can be obtained from a subgraph of $G$ by contracting edges.

In this paper we are concerned with the structure of the graphs which have no minor isomorphic to a given graph, and we begin in this section with a discussion of what kind of structure we might expect. It turns out that there are (at least) four ingredients to ``structure'', which we now review.

\textit{Ingredient} 1: \textit{Tree-structure.} It is a triviality that the graphs with no minor isomorphic to a loop are the forests; this is the simplest instance of the exclusion of a minor forcing some variety of tree-structure. If we exclude $K_3$ instead we get almost the same answer---forests augmented by loops and multiple edges. Excluding $K_4$%
\origfootnote[1]{Research partially supported by NSF Grant No. DMS 8504054.}%
\origfootnote[2]{Current address: Princeton University, Princeton, NJ 08544, USA.}
